\documentclass[11pt,a4paper]{article}

\usepackage[T1]{fontenc}
\usepackage[margin=2.4cm,top=2.5cm,bottom=2.5cm]{geometry}
\usepackage{newtxtext,newtxmath}
\usepackage[scaled=0.85]{inconsolata}
\usepackage{microtype}
\usepackage{booktabs}
\usepackage{tabularx}
\usepackage{array}
\usepackage{amsmath}
\usepackage{xcolor}
\usepackage{enumitem}
\usepackage{titlesec}
\usepackage{caption}
\usepackage{listings}
\usepackage{fancyhdr}
\usepackage{underscore}
\usepackage[colorlinks=true,linkcolor=accent,urlcolor=accent]{hyperref}

\definecolor{accent}{HTML}{1F6F6B}
\definecolor{rust}{HTML}{A8442A}
\definecolor{ink}{HTML}{16191A}
\definecolor{mute}{HTML}{5E6A6B}
\definecolor{pale}{HTML}{E8F0EF}
\definecolor{rule}{HTML}{C8D4D3}

\titleformat{\section}{\normalfont\large\bfseries\color{accent}}{\thesection}{0.7em}{}
\titleformat{\subsection}{\normalfont\normalsize\bfseries\color{ink}}{\thesubsection}{0.6em}{}
\titlespacing*{\section}{0pt}{1.5em plus .3em minus .2em}{0.55em}
\captionsetup{font=small,labelfont={bf,color=accent}}

\pagestyle{fancy}\fancyhf{}
\renewcommand{\headrulewidth}{0.4pt}
\fancyhead[L]{\small\color{mute}What \texttt{eACBN0.f90} does}
\fancyhead[R]{\small\color{mute}\thepage}

\lstset{
  basicstyle=\ttfamily\scriptsize, breaklines=true, columns=fullflexible,
  keepspaces=true, backgroundcolor=\color{pale!60}, frame=leftline,
  framerule=1.2pt, rulecolor=\color{accent}, xleftmargin=8pt,
  framexleftmargin=8pt, aboveskip=8pt, belowskip=8pt,
  commentstyle=\color{mute}\itshape,
}
\newcommand{\C}[1]{\texttt{\small #1}}

\begin{document}

\begin{center}
{\LARGE\bfseries What \texttt{eACBN0.f90} actually does}\\[4pt]
{\large\color{mute} The ACBN0 kernels inside \texttt{hp.x}, and what they\\
re-implement from \texttt{PW/src/new\_ns\_v.f90}}\\[9pt]
{\small Code report \quad\textbullet\quad 14 September 2026}
\end{center}

\vspace{2pt}
\noindent\textcolor{rule}{\rule{\textwidth}{0.8pt}}
\vspace{5pt}

\noindent
This is a description of the code, not of the validation. It answers two
questions: what the kernels compute, and which parts of them are transplanted
from the fork versus written here. Method reference: S.~Lee and Y.-W.~Son,
\emph{Phys.\ Rev.\ Research} \textbf{2}, 043410 (2020); ACBN0 itself: Agapito,
Curtarolo and Buongiorno Nardelli, \emph{Phys.\ Rev.\ X} \textbf{5}, 011006
(2015).

\vspace{2pt}
\noindent\textcolor{rule}{\rule{\textwidth}{0.8pt}}

% =====================================================================
\section{The quantity being computed}

ACBN0 is not linear response. It never perturbs anything. It evaluates a
\emph{ratio} on a converged ground state: the Coulomb (and exchange) integral
over the Hubbard manifold, weighted by \emph{renormalised} occupations, divided
by the same integral weighted by \emph{bare} occupations. The numerator is
screened by how much of each Bloch state actually lives on the manifold; the
denominator is not. Their ratio is the effective $U$.

\subsection{The working formulae, as the code has them}

Write $P^\sigma_{m m'}$ for the renormalised on-site occupation matrix,
$n^\sigma_{mm}$ for the bare one, and $(m_3 m_1 | m_4 m_2)$ for a two-electron
repulsion integral over real spherical harmonics. Then for a manifold that is
not a bare $s$ shell:
%
\begin{align}
A_0^U &= \sum_{m_1 m_2 m_3 m_4}
  \Big[P^{\uparrow}_{m_3m_1}P^{\uparrow}_{m_4m_2}
     + P^{\downarrow}_{m_3m_1}P^{\downarrow}_{m_4m_2}
     + P^{\uparrow}_{m_3m_1}P^{\downarrow}_{m_4m_2}
     + P^{\downarrow}_{m_3m_1}P^{\uparrow}_{m_4m_2}\Big]\,(m_3m_1|m_4m_2)
     \label{eq:a0u}\\[2pt]
A_0^J &= \sum_{m_1 m_2 m_3 m_4}
  \Big[P^{\uparrow}_{m_3m_1}P^{\uparrow}_{m_4m_2}
     + P^{\downarrow}_{m_3m_1}P^{\downarrow}_{m_4m_2}\Big]\,(m_3m_2|m_4m_1)
     \label{eq:a0j}
\end{align}
%
and the denominators, built from bare occupations only,
%
\begin{align}
\mathrm{HF}^U &= \sum_{m_1 \neq m_2}\!\big[n^{\uparrow}_{m_1}n^{\uparrow}_{m_2}
  + n^{\downarrow}_{m_1}n^{\downarrow}_{m_2}
  + n^{\uparrow}_{m_1}n^{\downarrow}_{m_2}
  + n^{\downarrow}_{m_1}n^{\uparrow}_{m_2}\big]
  \;+ \sum_{m_1 = m_2}\!\big[n^{\uparrow}_{m_1}n^{\downarrow}_{m_2}
  + n^{\downarrow}_{m_1}n^{\uparrow}_{m_2}\big] \\[2pt]
\mathrm{HF}^J &= \sum_{m_1 \neq m_2}\big[n^{\uparrow}_{m_1}n^{\uparrow}_{m_2}
  + n^{\downarrow}_{m_1}n^{\downarrow}_{m_2}\big]
\end{align}
%
giving
\[
\boxed{\;U \;=\; \frac{A_0^U}{\mathrm{HF}^U} \;-\; \frac{A_0^J}{\mathrm{HF}^J}\;}
\qquad\text{and, for a pure $s$ manifold,}\qquad U = \frac{A_0^U}{\mathrm{HF}^U}.
\]

\noindent
Two structural points fall straight out of these expressions.

\begin{itemize}[leftmargin=1.3em,itemsep=3pt,topsep=4pt]
\item The $s$ case drops the exchange term \emph{and} the same-spin part of the
      direct term. A single $s$ orbital has no orbital exchange partner, and the
      same-spin Hartree contribution is cancelled by Pauli exclusion, so only
      the $\uparrow\downarrow$ products survive. In the code this is the
      \C{IF ( lch(m1,nt) == 1 .AND. lch(m2,nt) == 1 )} branch --- and note it
      tests \emph{both} indices, because with background channels $m_1$ and
      $m_2$ can sit in different $\ell$.
\item The exchange integral has its indices transposed relative to the direct
      one: $(m_3m_1|m_4m_2)$ against $(m_3m_2|m_4m_1)$. That single swap is the
      whole difference between $A_0^U$ and $A_0^J$ in the source.
\end{itemize}

\subsection{The inter-site term}

For $V$ the same construction runs over \emph{pairs} of atoms $(I,J)$ in
neighbouring cells. The occupation matrix gains three indices ---
\C{ns(m1,m2,spin,na1,na2,nnc)} instead of \C{ns(m1,m2,spin,na)} --- and the
two-electron integral is evaluated with its two charge centres at
$\tau_{I}$ and $\tau_{J}+\mathbf{R}$ rather than both at $\tau_I$. The
formula is simpler than the on-site one:
\[
V_{IJ} \;=\; \frac{1}{2}\,\frac{E_0^{V}(I,J)}{\mathrm{HF}^{V}(I,J)} .
\]
No exchange term, and a factor $\tfrac12$ because the pair is accumulated twice,
once as $(I,J)$ and once as $(J,I)$.

% =====================================================================
\section{How the $+U$ kernel runs}

\C{hp_acbn0_ns_a0} is about 400 lines. Its shape:

\begin{enumerate}[leftmargin=1.6em,itemsep=3pt]
\item \textbf{Project.} For each $k$-point, read the KS wavefunctions
      (\C{iunwfc}) and the ortho-atomic Hubbard projectors (\C{iunhub}), and
      form $\langle \varphi_m | \psi_{n\mathbf{k}}\rangle$ with \C{calbec}.
      This is the only place the wavefunctions are touched.
\item \textbf{Bare occupations $n$.} Accumulate
      $n^\sigma_{m_1m_2} = \sum_{n\mathbf{k}} w_{n\mathbf{k}}\,
      \langle\psi|\varphi_{m_2}\rangle\langle\varphi_{m_1}|\psi\rangle$.
      Standard DFT+U occupation matrix.
\item \textbf{Per-band manifold weight.} For every band separately, accumulate
      $|\langle\varphi_m|\psi_{n\mathbf{k}}\rangle|^2$ into
      \C{num(ibnd)\%O}, indexed by $\ell$ and by \emph{type} rather than by
      atom, then sum it over all atoms of the same species
      (\C{upf(nt1)\%psd == upf(nt2)\%psd}). This is the ACBN0 renormalisation
      factor: how much of band $n\mathbf{k}$ sits on that species' manifold.
\item \textbf{Renormalised occupations $P$.} Repeat step 2, but weight each band
      by the factor from step 3. That is the single substantive difference
      between $n$ and $P$, and the entire physical content of ACBN0.
\item \textbf{Symmetrise.} Both $n$ and $P$ go through \C{hp_acbn0_symm_n4}.
\item \textbf{Assemble.} Build $A_0$ from $P$ and HF from $n$ per
      Eqs.~(1)--(4), then take the ratio.
\end{enumerate}

\subsection{An observation about the renormalisation}

Step 3 fills \emph{only} the diagonal \C{num(ibnd)\%O(l,l,...)}. Step 4 then
reads \C{num(ibnd)\%O(l1,l2,...)} with $\ell_1$ from $m_1$ and $\ell_2$ from
$m_2$. So whenever $m_1$ and $m_2$ lie in \emph{different} channels, the
renormalisation weight is zero and the corresponding block of $P$ vanishes by
construction.

For every single-channel system this is invisible: $\ell_1 = \ell_2$ always. It
only bites on a multi-channel manifold such as GaAs (Ga $4s+4p$), where it means
the $s$--$p$ cross-block of $P$ is identically zero while the $s$--$p$ block of
the \emph{bare} occupations is not. \textbf{This is the original code's
behaviour and we reproduce it exactly} --- which is precisely why our GaAs run
matches \C{pw.x} across all 204 lines. I am flagging it as a property of the
method as implemented, not as a defect, because I have not checked it against
the ACBN0 paper's intent and it is the kind of thing worth confirming with the
author before anything is published on multi-channel manifolds.

\subsection{The two-electron integrals}

The integrals come from outside this file. \C{pao_3g} returns a three-Gaussian
fit to the pseudo-atomic orbital for a given species and $\ell$, and
\C{compute_eri} evaluates $(ab|cd)$ analytically over those Gaussians at four
given centres. \C{hp_acbn0_orberi} is the small bookkeeping routine that walks
the real-spherical-harmonic index: $1$ for $s$; $2,3,4$ for $p$; $5\ldots 9$ for
$d$. It is a stateful cycler --- called once per $m$ inside each loop, it
advances its argument, so the $m$-th orbital of the manifold gets the $m$-th
harmonic index.

This is also where the $s,p,d$-only restriction comes from: \C{pao_3g} has no
$f$-shell fit and returns zero coefficients \emph{silently}, which would make
every integral vanish and $U$ come out as $0/0$. Both check routines refuse
$\ell > 2$ explicitly for that reason.

% =====================================================================
\section{Provenance: what is transplanted, what is new}

\subsection{Transplanted from \texttt{PW/src/new\_ns\_v.f90}}

\begin{center}\small
\begin{tabular}{llrl}
\toprule
Our routine & From & Stmts & Role \\
\midrule
\C{hp_acbn0_ns_a0}    & \C{new_ns_a0} 233--635        & 176 & the $+U$ kernel \\
\C{hp_acbn0_ns_e0_1}  & \C{new_ns_e0_1} 636--1808     & 577 & the $+U+V$ kernel \\
\C{hp_acbn0_ns_v}     & \C{new_ns_v} 10--232          & 103 & occupation matrices alone \\
\C{hp_acbn0_orberi}   & \C{orberi} 3780--3819         &  25 & harmonic index cycler \\
\C{hp_acbn0_symm_n4}  & \C{symmetrize_n4} 2964--3167  &  86 & symmetrise $n(m_1,m_2,\sigma,I)$ \\
\C{hp_acbn0_symm_n4m} & \C{symmetrize_n4m} 3168--3399 & 104 & ditto, magnetic \\
\C{hp_acbn0_symm_n5}  & \C{symmetrize_n5} 3400--3779  & 174 & symmetrise the 6-index pair matrix \\
\midrule
\multicolumn{2}{l}{\textbf{total transplanted}} & \textbf{1245} & \\
\bottomrule
\end{tabular}
\end{center}

\noindent
\textbf{How faithful.} The transplanted bodies were diffed against the fork by a
script (\C{normdiff4.py}) that joins Fortran continuations into whole
statements, strips comments, keeps only executable code, collapses the
mechanical renames to common tokens, and canonicalises the intended systematic
edits. \textbf{1245 statements, 18 differing}, in three classes:

\begin{itemize}[leftmargin=1.3em,itemsep=2pt,topsep=3pt]
\item 6 --- \C{hp_acbn0_orberi}'s dummy argument renamed \C{ehub_l}
      $\rightarrow$ \C{lp1} (cosmetic; it is a dummy name).
\item 8 --- the dropped \C{stop_acbn0} convergence short-circuit and its
      \C{ehub_v_diff} bookkeeping. \C{hp.x} is one-shot and has no SCF history,
      so there is nothing to short-circuit.
\item 4 --- the two bug-fix guards (\S\ref{sec:bugs}).
\end{itemize}

\noindent
\C{hp_acbn0_ns_a0} and \C{hp_acbn0_ns_v} come out at \textbf{zero}
differences. The measurement is stronger than it looks: since stage V4 the
canonicaliser no longer collapses the $m$ index, so every
\C{ehub_l(m1,nt)} in the fork must match \C{lch(m1,nt)} here literally, index by
index --- all 225 of them do.

\subsection{Written here}

\begin{center}\small
\begin{tabularx}{\textwidth}{lX}
\toprule
\C{MODULE hp_acbn0_nn} & The neighbour table. Reproduces \C{hubbard_nns()} and
the \C{nnnc} count from \C{input.f90:2144}, but built from \C{tau}/\C{at}/\C{alat}
and a cutoff rather than read from \C{ldaV}. Also \C{hp_acbn0_symmetry} and
\C{hp_acbn0_get_nnc}, from \C{PW/src/hubbard_nns.f90}. \\
\addlinespace
\C{MODULE hp_acbn0_man} & The manifold. Maps official \C{ldaU}'s \C{ll(m,nt)},
\C{ldim_u}, \C{ldmx_tot} and the three projector blocks
(\C{offsetU}/\C{_back}/\C{_back1}) onto the contiguous view
\C{lm}/\C{lch}/\C{ofs} that the kernels are written against. Entirely new ---
\C{ldaV} stored the manifold contiguously, so the fork never needed it. \\
\addlinespace
drivers & \C{hp_acbn0}, \C{hp_acbn0_u}, \C{hp_acbn0_v}: buffer handling,
\C{orthoUwfc}, \C{d_matrix}, allocation, cleanup. Replaces what \C{sum_band}
did for the fork. \\
\addlinespace
guards & \C{hp_acbn0_check}, \C{hp_acbn0_check_v}: preconditions the kernels
assume silently (ortho-atomic projectors, $n_{\rm spin}=2$, $s/p/d$ only,
no \C{gamma_only}). Each of these fails as a \emph{wrong number} rather than a
crash if unchecked. \\
\addlinespace
output & \C{hp_acbn0_write}, \C{_write_v}, \C{_print_uv}, \C{hp_acbn0_manlab}.
\C{_print_uv} deliberately reproduces \C{pw.x}'s own table format so that every
comparison is a plain \C{diff}. \\
\bottomrule
\end{tabularx}
\end{center}

\subsection{Deliberately not transplanted}

\C{new_ns_e0_2} (\C{acbn0_type = 3}) and the entire noncollinear set ---
\C{new_ns_v_nc}, \C{new_ns_e0_nc_1}, \C{symmetrize_n4_nc}, \C{_n4m_nc},
\C{_n5_nc}, roughly 2300 further lines. This is why the authors' GaAs reference,
which is noncollinear with spin--orbit, remains unreachable.

% =====================================================================
\section{Four bugs found in the original}\label{sec:bugs}

The fork's patch touches 60 files and \textbf{none is under \C{HP/}} --- ACBN0
was implemented in \C{pw.x} only. \C{pw.x} initialises \C{MODULE ldaV} from two
paths (the input file, and the XML); \C{hp.x} only ever walks the second. Every
bug below is a consequence.

\begin{description}[leftmargin=0pt,itemsep=5pt]
\item[\textcolor{rust}{Bug 1 --- \C{remove_ehub_u} / \C{_v} undefined (silent).}]
No initialiser in \C{ldaV.f90}; assigned only at \C{input.f90:2083}. The kernels
multiply the final $U$ by them \emph{twice}. Undefined, this does not crash: it
silently scales or zeroes every Hubbard parameter. Not hypothetical --- the
authors' own \C{GaAs.acbn0.in_noncollienar} sets \C{remove_ehub_u(1,1) = 0}.

\item[\textcolor{rust}{Bug 2 --- \C{iunwfc} closed at the insertion point (crash).}]
\C{read_file} buffers the wavefunctions then closes the unit with
\C{close_buffer(...,'KEEP')}. \C{'KEEP'} preserves the \emph{files}, but the call
also deregisters the unit, so \C{get_buffer} falls through to \C{davcio} on a
closed unit. Official \C{hp.x} never trips this because DFPT reads through
\C{iuwfc}, opened later.

\item[\textcolor{rust}{Bug 3 --- \C{hubbard_get_nnc} returns an undefined index.}]
\C{hubbard_nns.f90:241}. \C{nnc} is assigned only inside
\C{IF (MAXVAL(...) == 1)}, so a cell carrying no neighbour pair falls through
untouched. \C{hubbard_symmetry} then asks for the index of a \emph{rotated}
cell, which need not be occupied, and \C{symmetrize_n5} indexes \C{n(nnc)} with
whatever was on the stack. Silent in an optimised build whenever the garbage
index lands in a mapped page; \C{-check bounds,uninit} named it at once.

\item[\textcolor{rust}{Bug 4 --- \C{l_separation} writes one element before the array (latent, still live).}]
In \C{symmetrize_n4}, \C{_n4m} and \C{_n5}, a type with no Hubbard manifold has
\C{ehub_l(1,nt) = 0} and the next line writes \C{l_separation(0,nt)}. Dormant
when every type is Hubbard-active --- true of NiO, MnO, ZnO, CdO, CrI$_3$ and
GaAs --- and it fires the moment a spectator species is added.
\end{description}

\noindent
Bugs 3 and 4 are the two guards that show up in the static diff. Bugs 1 and 2
needed no guard: bug 1 disappears because \C{ldaU} has no such mask, and bug 2
is handled by reopening the unit in the driver.

% =====================================================================
\section{Where to look in the file}

\begin{center}\small
\begin{tabular}{rll}
\toprule
Line & Routine & \\
\midrule
   1 & file header & the \C{ldaV} $\rightarrow$ \C{ldaU} translation table \\
  94 & \C{MODULE hp_acbn0_nn} & neighbour table \\
 491 & \C{MODULE hp_acbn0_man} & manifold, and the \C{wfcU} column map \\
 704 & \C{hp_acbn0_u} & $+U$ driver \\
 790 & \C{hp_acbn0_check} & $+U$ preconditions \\
 890 & \C{hp_acbn0_write} & \C{eACBN0_hub.dat} \\
1031 & \C{hp_acbn0_ns_a0} & \textbf{the $+U$ kernel} \\
1436 & \C{hp_acbn0_symm_n4} & symmetrisation \\
1648 & \C{hp_acbn0_orberi} & harmonic index cycler \\
1701 & \C{hp_acbn0} & dispatcher on the \C{acbn0} tag \\
1721 & \C{hp_acbn0_v} & $+U+V$ driver \\
2215 & \C{hp_acbn0_ns_e0_1} & \textbf{the $+U+V$ kernel} \\
3372 & \C{hp_acbn0_ns_v} & occupation matrices alone \\
3596 & \C{hp_acbn0_symm_n4m}, 3839 \C{_n5} & symmetrisation \\
\bottomrule
\end{tabular}
\end{center}

\noindent
4234 lines, two modules, 22 subroutines and one function. Seven other files in
\C{HP/src/} carry small edits: \C{hpcom.f90} (two variables),
\C{hp_readin.f90} (namelist, defaults, four relaxed DFPT guards),
\C{hp_bcast_input.f90} (two broadcasts), \C{hp_main.f90} (the dispatch, placed
above \C{hp_init}), \C{hp_print_clock.f90}, \C{Makefile}, \C{make.depend}.

% =====================================================================
\section{Two things I would change}

\noindent
\textbf{Stale guard exemptions in \C{hp_readin.f90}.} Three guards are written
\C{.AND. acbn0 /= 2}, which was correct at stage V2 when the $+V$ branch still
ran on \C{ldaV}. Since V3 it runs on \C{ldaU} and \emph{does} need
\C{lda_plus_u}, ortho-atomic projectors and \C{lda_plus_u_kind} $\neq 1$. Not a
correctness bug --- all three are backstopped by \C{hp_acbn0_check_v}, so a bad
input gets a clear error rather than a wrong number --- but the exemptions are
dead weight and the comments are two stages out of date.

\smallskip
\noindent
\textbf{\C{remove_ehub_u} has no route in.} It is input-only and is not
round-tripped through the XML, so we omit the factors entirely --- exactly
equivalent whenever the \C{pw.x} run used the documented default of 1, and
wrong for any run that masks a channel, as the authors' GaAs input does. The fix
is one \C{\&INPUTHP} tag.

\vspace{8pt}
\noindent\textcolor{rule}{\rule{\textwidth}{0.8pt}}
\vspace{2pt}

\noindent\footnotesize\color{mute}
Companion documents: \C{REPORT-eacbn0-hpx.pdf} (validation), \C{slides.html}
(group meeting), \C{logs/260901-eacbn0-hpx.txt} and \C{logs/260904-eacbn0-v4.txt}
(session logs). The faithfulness measurement is re-runnable:
\C{python3 normdiff4.py}.

\end{document}
